% Verbatim from /Users/alingalatan/Downloads/cubic_chaos_partial_results (2).zip:direct3_explicit_local.tex (16 November 2026); GPT (a ChatGPT session of the user), relayed by the user.
% /workspace/scratch/3efc7fe714fa/direct3_explicit_local.tex
\documentclass[11pt]{article}
\usepackage{amsmath,amssymb,amsthm,hyperref}
\newcommand{\E}{\mathbb E}\newcommand{\Pp}{\mathbb P}\newcommand{\Ga}{\Gamma}
\newtheorem{theorem}{Theorem}\newtheorem{lemma}{Lemma}
\begin{document}
\title{An explicit local cubic-chaos $L^1$ bound by localized integration by parts}
\author{}\date{}\maketitle
All statements below are rigorous; constants are deliberately very loose.
Write $c=\sqrt{2/\pi}$. The normalization of multiple integrals is
$\E I_q(f)^2=q!\|f\|^2$.

\begin{theorem}
For $F\in H_3$, $\E F^2=1$, put $d=\E F^4-3$. If $0<d\le1$, then
\[
 \E|F|\le c(1-d/24)+10^{60}d^{5/4}.
\]
In particular, for $0<d\le10^{-248}$,
\[
 \E|F|\le c-\frac{c}{48}d<c.
\]
\end{theorem}

\paragraph{Sources and conventions.}
We use only the usual Gaussian integration by parts, hypercontractivity,
normal Stein derivative bounds (equations (1.4)--(1.5) of Gaunt,
\emph{On Stein factors in Stein's method for normal approximation},
\url{https://arxiv.org/html/2410.05475v1}), and the explicit contraction inequality
(4.50) in Bierme--Bonami--Nourdin--Peccati,
\emph{Optimal Berry--Esseen rates on the Wiener space: the barrier of third and fourth cumulants},
\url{https://arxiv.org/abs/1109.1546}.
Its formula (3.32)/(4.47) is used to count the coefficients below.
The bound $\operatorname{Var}(\|DF\|^2/3)\le 2d/9$ and the TV bound
$d_{\rm TV}(F,G)\le2\sqrt{\operatorname{Var}(\|DF\|^2/3)}$
appear in (1.1)--(1.2) of that paper. Here TV distance is $\sup_B|\mu(B)-\nu(B)|$.
The nonsmooth test $|x|$ is never inserted directly into a higher derivative identity.

\paragraph{1. Gamma operators and an explicit contraction constant.}
Set $\Ga_0=F$ and
$\Ga_{j+1}=\langle DF,-DL^{-1}\Ga_j\rangle$, with $L^{-1}$ annihilating constants.
Then $A:=\Ga_1=\|DF\|^2/3$, $\E A=1$,
$\E\Ga_2=\kappa_3/2=0$, $\E\Ga_3=d/6$.
By Gaussian parity, $\kappa_3=0$. For a smooth function $f$ of bounded derivatives,
\begin{equation}\label{expansion}
\E[f'(F)-Ff(F)]= -\frac d6\E f'''(F)-\E[\Ga_4 f''''(F)].
\end{equation}
For $F=I_3(t)$, formula (4.47) expresses $\Ga_4$ as a sum of at most $3^4=81$
multiple integrals, each of order at most $7$. Every coefficient is at most
$12\cdot60^3$: the initial coefficient is at most $12$, and each later multiplier
is $3(r-1)!\binom2{r-1}\binom{k-1}{r-1}$ with $r\le3$ and $k-1\le5$,
hence at most $60$. Inequality (4.50) bounds each kernel norm by
$\max_{r=1,2}\|t\otimes_rt\|^{5/2}\le d^{5/4}$ (we discard the stronger constant in (4.43)).
The integral isometry and $\sqrt{7!}<71$ therefore give
\begin{equation}\label{gamma_bound}
 \|\Ga_4\|_2\le81\cdot12\cdot60^3\cdot71\,d^{5/4}<10^{11}d^{5/4}.
\end{equation}
Also $\Ga_4$ is a polynomial of degree at most $7$.

\paragraph{2. Quantitative localized integration by parts.}
Define $\chi(a)=0$ for $a\le1/2$, $\chi(a)=1$ for $a\ge3/4$, and
$\chi(a)=3u^2-2u^3$ for $u=4a-2\in[0,1]$. Thus
$0\le\chi\le1$, $|\chi'|\le6$, $|\chi''|\le96$ in the weak sense.
Define, for $a\ge0$,
\[
 v(a)=\begin{cases}2-4(a-1/2)+8(a-1/2)^2,&a\le1/2,\\1/a,&a\ge1/2.\end{cases}
\]
This is $C^2$, with $|v|\le6$, $|v'|\le12$, $|v''|\le16$.
Put
\[
 V=\tfrac13v(A)DF,\quad J=\langle DA,DF\rangle,\quad
 B=v(A)F-\tfrac13v'(A)J,\quad TZ=BZ-\langle V,DZ\rangle.
\]
Here $T Z=\delta_{\rm Mall}(ZV)$, where $\delta_{\rm Mall}$ denotes Gaussian divergence.
For $Z$ supported on $\{A\ge1/2\}$,
\begin{equation}\label{ibp}
 \E[Z\psi'(F)]=\E[(TZ)\psi(F)],\qquad
 \E[Z\psi''(F)]=\E[(T^2Z)\psi(F)].
\end{equation}
Indeed $\langle DF,V\rangle=Av(A)=1$ on this support, and $TZ$ remains supported there.
The $W^{2,\infty}$ cutoff is legitimate by Sobolev integration by parts, or by smooth approximation.

Here are explicit, dimension-free bounds sufficient for both uses of \eqref{ibp}:
\begin{equation}\label{Tbound}
 \E|T^2\chi(A)|\le10^{32},\qquad
 \E|T^2(\Ga_4\chi(A))|\le10^{32}\|\Ga_4\|_2.
\end{equation}
For completeness, the numerical bookkeeping follows. For $p\le32$, $j\le2$,
\[
 \|D^jF\|_p\le10^4,\qquad \|D^jA\|_p\le10^4.
\]
This follows from hypercontractivity and
$\|A\|_2^2=1+\operatorname{Var}A\le1+2/9<4$;
for a degree-$m$ polynomial $P$,
$\|D^jP\|_p\le m^{j/2}(p-1)^{m/2}\|P\|_2$.
The same argument applies to Hilbert-valued derivatives, by Minkowski's inequality.
With $M=10^4$, Holder's inequality gives
\[
 \|V\|_8\le2M,\quad \|DV\|_8\le5\cdot10^8,\quad
 \|B\|_8\le10^9,\quad \|DB\|_8\le10^{13}.
\]
For example,
$\|DB\|_8\le6M+20M^2+(16/3)M^3<10^{13}$;
products of three factors need at most their $L^{24}$ norms.
Expanding $T^2$ gives
\[
 T^2 Z=(B^2-\langle V,DB\rangle)Z-2B\langle V,DZ\rangle
 +\langle DZ,(DV)V\rangle+D^2Z[V,V].
\]
Consequently
\[
 \E|T^2Z|\le2\cdot10^{18}
 (\|Z\|_2+\|DZ\|_2+\|D^2Z\|_2).
\]
For $Z=\chi(A)$ the sum is below $2\cdot10^{10}$.
For $Z=\Ga_4\chi(A)$, writing $Q=\|\Ga_4\|_2$, the product rule gives
\[
 \|Z\|_2\le Q,\qquad \|DZ\|_2\le3\cdot10^6Q,\qquad
 \|D^2Z\|_2\le3\cdot10^{12}Q.
\]
In the last estimate use
$\|\Ga_4\|_4<47Q$, $\|\Ga_4\|_6<280Q$,
$\|D\Ga_4\|_4<124Q$, and $\|D^2\Ga_4\|_2\le7Q$.
Thus \eqref{Tbound} follows, with considerable room.

\paragraph{3. A localized density estimate.}
Hypercontractivity for the degree-four polynomial $A-1$ gives
\begin{equation}\label{bad}
 p_*:=\Pp(A<3/4)\le4^{16}\|A-1\|_{16}^{16}
 \le900^{16}d^8<10^{48}d^8.
\end{equation}
We shall also use $p_*\le16\operatorname{Var}A<4d$.
Let $\mu_\chi(E)=\E[\chi(A)1_{\{F\in E\}}]$.
Two integrations by parts show that it has a Lipschitz density $p_\chi$ with
\[
 p_\chi(x)=\E[1_{\{F>x\}}T\chi(A)],\qquad
 p_\chi'(x)=-\E[1_{\{F>x\}}T^2\chi(A)]
\]
in the weak/a.e. sense; hence $\operatorname{Lip}(p_\chi)\le10^{32}$.
Writing $\phi$ for the standard normal density,
\[
 \|p_\chi-\phi\|_1\le2d_{\rm TV}(F,G)+p_*
 \le2\sqrt d+4d\le6\sqrt d.
\]
If an integrable function has Lipschitz constant $L$, its supremum norm is at most
$\sqrt{L\|u\|_1}$, by integrating the triangular lower bound around a point
where $|u|$ nearly attains its supremum. Here $\operatorname{Lip}(p_\chi-\phi)\le2\cdot10^{32}$. It follows that
\begin{equation}\label{density}
 \|p_\chi-\phi\|_\infty\le4\cdot10^{16}d^{1/4}.
\end{equation}
No boundedness or smoothness of the density of the discarded portion of $F$ is assumed.

\paragraph{4. Apply Stein to a smoothed absolute value.}
Let $e=d^2$, $h_e(x)=\E|x+eG|$, and let $f_e$ be the usual normal Stein solution:
\[
 f_e'(x)-xf_e(x)=h_e(x)-\E h_e(G).
\]
The standard Stein derivative estimates for smooth Lipschitz tests give
\[
 \|f_e'\|_\infty\le2,\quad \|f_e''\|_\infty\le2,\quad
 \|f_e''''\|_\infty\le4e^{-2}.
\]
For the last inequality, use $\|f_h^{(k+1)}\|_\infty\le2\|h^{(k)}\|_\infty$
with $k=3$, and $h_e'''=2\phi_e'$, where $\phi_e$ is the $N(0,e^2)$ density.
From \eqref{Tbound}, \eqref{bad}, and \eqref{gamma_bound},
\begin{align*}
 |\E\Ga_4 f_e''''(F)|
 &\le2\E|T^2(\Ga_4\chi(A))|
       +4e^{-2}\|\Ga_4\|_2\sqrt{p_*}\\
 &\le (2\cdot10^{32}+4\cdot10^{24})10^{11}d^{5/4}
 <10^{44}d^{5/4}.
\end{align*}

It remains to compare $\E f_e'''(F)$ and $\E f_e'''(G)$ without a derivative loss.
Differentiating the Stein equation twice gives
\[
 f_e'''(x)=2\phi_e(x)+x f_e''(x)+2f_e'(x).
\]
By \eqref{density} and \eqref{bad},
\[
 |\E\phi_e(F)-\E\phi_e(G)|
 \le4\cdot10^{16}d^{1/4}+p_*/(\sqrt{2\pi}e)
 \le4\cdot10^{16}d^{1/4}+10^{48}d^6
 <2\cdot10^{48}d^{1/4}.
\]
For the remaining terms, the full variation norm of $\mathcal L(F)-\mathcal L(G)$
is at most $2\sqrt d$. Cauchy--Schwarz against its total variation measure and
$\E F^2+\E G^2=2$ yield
\[
 |\E[Ff_e''(F)]-\E[Gf_e''(G)]|\le4d^{1/4},\qquad
 2|\E f_e'(F)-\E f_e'(G)|\le8\sqrt d.
\]
Thus, very conservatively,
\begin{equation}\label{third}
 |\E f_e'''(F)-\E f_e'''(G)|<10^{49}d^{1/4}.
\end{equation}
The Gaussian coefficient is exact:
\[
 \E f_e'''(G)=-\tfrac14\E[h_e(G)H_4(G)]
 =\frac{c}{4(1+e^2)^{3/2}}.
\]
Since $|x|\le h_e(x)$, \eqref{expansion} and \eqref{third} give
\[
 \E|F|\le c\sqrt{1+e^2}
 -\frac{cd}{24(1+e^2)^{3/2}}
 +\left(10^{49}/6+10^{44}\right)d^{5/4}.
\]
Using $e=d^2$, $\sqrt{1+e^2}\le1+e^2/2$, and
$(1+e^2)^{-3/2}\ge1-3e^2/2$, the first two terms are at most
$c(1-d/24)+(1/2)d^4+(1/16)d^5$.
This proves the claimed $10^{60}$ bound for $0<d\le1$.
Finally $10^{60}(10^{-248})^{1/4}=10^{-2}<c/48$, proving the stated strict inequality.

\paragraph{Scope.}
This is only a local theorem. It does not settle the range $d>10^{-248}$.
For arbitrary $h$, apply it to $F=h/\|h\|_2$ and use the standardized excess fourth cumulant.
The $d=1$ Gaussian check gives equality and no contradiction: its fourth cumulant is zero.
For $H_3(G)/\sqrt6$ and products of three Gaussians the standardized cumulants are not in this tiny range;
no conclusion from the local threshold is claimed for those examples.
\end{document}
